\section{The Wider Fátima Cycle: First Saturdays and Consecration}

The July 1917 secret says that the Lady will return to ask for reparatory Communion on First Saturdays and the consecration of Russia. The received narrative locates those developments at Pontevedra and Tuy after Lúcia had left Fátima. They belong to the message's historical unfolding and have received substantial ecclesial reception. They do not thereby become part of the exact 1930 operative clause.

\subsection{The reported seventh encounter, 1921}

Before leaving for the Asilo de Vilar in Porto, Lúcia struggled with separation from family and the apparition places. Her later diary describes an encounter at the Cova on 15 June 1921, understood as the promised seventh visit: she was told to follow the path set by the bishop because it was God's will. Within the received narrative, spectacle yields to ecclesial obedience. The visionary's vocation is not to remain the proprietor of a holy place but to accept formation, obscurity, and authority.

The biographical report is outside the 1930 six-event clause. Its theological criterion remains public: obedience to a legitimate command binds within its competence and moral limits; no apparition can authorize abuse, falsehood, or blind submission. Here the direction to leave protected a young witness from attention and prepared her long ecclesial service.

\subsection{Pontevedra, 10 December 1925}

While Lúcia was a postulant with the Sisters of Saint Dorothy at Pontevedra, she reported seeing Mary and the Child Jesus. The narrative presents Mary's heart encircled by thorns and asks a five-month practice of reparation. A 1927 letter is the early written witness used by the Shrine.

\begin{fourstable}{Element}{Reported request}{Catholic meaning}{Necessary safeguard}
Five consecutive first Saturdays & Perseverance across a bounded sequence & Habitual conversion rather than a single emotional act & Missing a day is not divine rejection; the devotion may simply be begun again.\\
Confession & Sacramental reconciliation associated with the practice & Contrition, absolution, and restored communion with the Church & The sacrament follows Church law; private revelation does not alter validity or seal.\\
Holy Communion & Receive Christ with reparatory intention & The Eucharist is source and summit; reparation is participation in Christ & A person conscious of grave sin must not receive without prior sacramental confession except as law provides.\\
Five decades of the Rosary & Contemplate Christ's mysteries with Mary & Scriptural remembrance, praise, petition, and discipleship & Not a magic count or replacement for Mass.\\
Fifteen minutes of meditation & Keep Mary company while considering Rosary mysteries & Unhurried contemplative attention forms the heart & The historic reference to fifteen mysteries reflects the then-received Rosary; it does not prohibit meditation within the current twenty-mystery cycle.\\
Reparatory intention & Perform the acts in love for Mary's wounded heart & Repentance and love oppose sin's contempt and ingratitude & Intention cannot turn a sacrament into a transaction or pain into currency.\\
Promise of final assistance & Mary's maternal assistance with graces needed for salvation & Hope in intercession and perseverance within Christ's grace & No automatic salvation without faith, repentance, charity, or final perseverance.\\
\end{fourstable}

The received 1926 clarification allows confession within a nearby interval when the communicant is in grace and retains or renews the reparatory intention. That devotional accommodation cannot dispense from the divine and ecclesial conditions for worthy Communion. If mortal sin is present, confession is not merely a checklist item that can be postponed at will.

The promise is frequently distorted. Catholic promises attached to a devotion do not compel God by performance. Final perseverance is a gift lived through freedom; Mary assists as mother and intercessor wholly in Christ. The practice is authentic when it deepens sacramental conversion, not when it becomes a guarantee used to postpone repentance.

The reported request did not create the Eucharist, Marian veneration, or reparation; it joined inherited Catholic elements in a specific five-month form. Marian observance on Saturday is documented from the Carolingian period, although the historical reason for choosing Saturday remains uncertain. The Church's emph{Directory on Popular Piety and the Liturgy} therefore receives the First Five Saturdays as a pious practice while refusing a merely temporal or mechanical reading: Communion must be lived within the Eucharist and the Paschal Mystery. Private revelation may deepen and spread devotion, but it remains ordered to the sacramental economy already given in Christ.

\subsection{Tuy, 13 June 1929}

In the chapel at Tuy, Lúcia reported a Trinitarian and Eucharistic vision. A cross of light filled the space. Its figures represented the Father, the Holy Spirit under the sign of a dove, and the crucified Son. A host and chalice received blood from the Crucified. Mary stood beneath one arm with her Immaculate Heart; beneath the other, words signified “Grace and Mercy.” The Lady then reportedly said that the time had come for the Pope, united with the world's bishops, to consecrate Russia to her Immaculate Heart and promised salvation by this means.

The imagery controls the request. Father, Son, and Spirit enclose it. Christ crucified and Eucharistic is its source. “Grace and Mercy” name its atmosphere. Mary stands beneath the Cross, not above the Trinity. Consecration to her is therefore an entrustment inside baptismal consecration to God, asking that her \latin{fiat} form a people for Christ.

The account was given to Lúcia's confessor and belongs to the later cycle. Its papal reception is strong; the 1930 decree's wording remains what it is.

\subsection{What Marian consecration means}

Only God is holy in himself and makes holy. Christ consecrates himself to the Father for his disciples so that they may be consecrated in truth (Jn 17:19). Baptism configures the person to Christ. No later devotional act creates a superior ownership or transfers a nation outside Christ's lordship.

Marian “consecration” is therefore analogical and relational: entrustment to Mary's maternal intercession, renewal of baptismal belonging, and adoption of her faithful consent. John Paul II's 1984 act makes the hierarchy explicit by uniting the Church's act with Christ's redemptive consecration. The Pope does not claim magical control over other persons' wills, invalidate legitimate civil sovereignty, or make Mary a source of sanctifying grace.

Consecrating a people also does not mean judging every member guilty of one ideology. It is intercession for persons and cultures, including persecutors and victims. The act asks grace to free, heal, reconcile, and order freedom toward truth.

\subsection{The major papal acts}

\begin{fourstable}{Date and pope}{Object and form}{Relation to Fátima}{Boundary}
31 Oct. 1942, Pius XII & Radio act entrusting and consecrating the Church and world to the Immaculate Heart & First major papal response during world war; renewed in the received history & Did not use every later-requested specification and should not be retrospectively misquoted.\\
7 July 1952, Pius XII & \emph{Sacro vergente anno} consecrates the peoples of Russia & Russia explicitly named within prayer for freedom and peace & Not presented by the CDF as the final answer to the collegial request.\\
7 June 1981 / 13 May 1982, John Paul II & Entrustments of the human family, developing after the assassination attempt & Papal reception becomes explicitly connected with Fátima and the third part & Preparatory stages, not evidence of disobedient bad faith.\\
25 Mar. 1984, John Paul II & World, all persons and peoples entrusted and consecrated in spiritual union with bishops convoked beforehand; nations in special need included & CDF 2000 publishes Lúcia's 1989 confirmation that it was done as requested & The public act need not satisfy a private commentator's preferred wording once the controlling dossier gives this answer.\\
25 Mar. 2022, Francis & Church and humanity, especially Russia and Ukraine, entrusted and consecrated during renewed war & A new collegial prayer answering a new historical wound & Renewal does not imply that the 1984 act was invalid or that devotion operates mechanically.\\
\end{fourstable}

The 1984 act did not speak “Russia” in the public text as some campaigners demanded, but it included those persons and nations especially awaiting entrustment and was performed in the convoked bishops' spiritual union. The CDF does not leave the matter suspended: it publishes Lúcia's written confirmation of 8 November 1989 and concludes that further discussion or request lacks basis. A Catholic may analyze the history; one may not represent the controlling ecclesial dossier as saying the opposite.

The 2022 act is best read as liturgical-pastoral repetition under a new war, just as baptismal promises and acts of entrustment can be renewed without declaring the first act void. Prayer enters history repeatedly because persons and nations remain free and wounded.

\subsection{Devotion without mechanism}

First Saturdays and consecrations have a common danger: they can be treated as techniques for controlling God. Their authentic form is the reverse. Confession surrenders self-justification; Communion receives the self-giving Christ; Rosary meditation submits imagination to the Gospel; reparation lets charity oppose sin; consecration yields autonomy to grace. The practices do not manipulate Providence. They dispose human freedom to cooperate with it.

No Catholic must complete the five Saturdays or use a particular consecration formula. Those who do should follow the Church's sacramental discipline, avoid scrupulosity, and let devotion bear visible fruit in worship, restitution, mercy, peace-making, and care for victims.
